%% ma: L11-B15-0005
%% lop: 11
%% chuong: 5
%% bai: 15
%% dang: 11.15.4
%% loai: TL
%% muc: VD
%% dap_an: "-"
%% nguon: "tai-lieu-hoc-tap-mon-toan-lop-11-theo-chuong-trinh-giao-duc-pho-thong-2018.pdf (Lớp Toán Thầy Duy, Phạm Lê Duy, THPT Chu Văn An), Chương 3 Bài 1, Dạng 3, Ví dụ 5 (tr.310)"
%% kien_thuc: "Lớp 11 Bài 15 (giới hạn của dãy số)"
%% trang_thai: cho-duyet
%% ly_do: "Dạng mới đề xuất (11.15.x), chờ giáo viên duyệt dạng; tài liệu gốc không có lời giải mẫu, đáp án do kho tự giải và kiểm bằng sympy."
\begin{ex}
Tính các giới hạn sau:

a) $\lim\dfrac{1+3^n}{3\cdot3^n+2^n}$;

b) $\lim\dfrac{4\cdot3^n-2^n}{2\cdot5^n-4^n}$;

c) $\lim\dfrac{7^n+1}{-2\cdot3^n-3\cdot6^n}$;

d) $\lim\dfrac{5^{n+1}-4^n+1}{-2\cdot5^n-6^n}$.
\loigiai{\sloppy a) Chia tử và mẫu cho $3^n$: $\dfrac{\left(\frac13\right)^n+1}{3+\left(\frac23\right)^n}\to\dfrac13$.

b) Chia cho $5^n$: \[\dfrac{4\left(\frac35\right)^n-\left(\frac25\right)^n}{2-\left(\frac45\right)^n}\to\dfrac01\cdot\dfrac12=0.\]

c) Chia cho $6^n$: $\dfrac{\left(\frac76\right)^n+\left(\frac16\right)^n}{-2\left(\frac12\right)^n-3}$; tử $\to+\infty$, mẫu $\to-3$ nên giới hạn bằng $-\infty$.

d) Chia cho $6^n$: \[\dfrac{5\left(\frac56\right)^n-\left(\frac23\right)^n+\left(\frac16\right)^n}{-2\left(\frac56\right)^n-1}\to\dfrac{0}{-1}=0.\]}
\end{ex}
